\chapter{激光冷却、Sisyphus 冷却、原子钟与磁阱}\label{ch:laser-cooling-clock}

\section{散射力、摩擦与反冲扩散}

冷却操纵的是原子质心自由度。下文用 $M$ 表示\emph{整颗原子}的质量；电子质量 $m_e$ 不出现在质心 Fokker--Planck 方程中。内部两能级稳态则直接继承\eqnrefs{eq:two-level-saturation}。

接下来考察两能级稳态散射力。

原子吸收一个沿 \(\mathbf k\) 方向的光子，动量改变

\[
\Delta\mathbf p=\hbar\mathbf k.
\]

自发辐射方向随机，长时间平均动量接近零。因此平均散射力

\[
\boxed{
\mathbf F_{sc}=\hbar\mathbf k\,R_{sc}
}
\]

其中散射率

\[
R_{sc}=\Gamma_{\rm sp}\rho_{ee}^{(\rm ss)}.
\]

于是

\[
\boxed{
\mathbf F_{sc}
=\hbar\mathbf k\frac{\Gamma_{\rm sp}}2
\frac{s_0}{1+s_0+[2(\Delta-\mathbf k\cdot\mathbf v)/\Gamma_{\rm sp}]^2}
}.
\]

Doppler shift 通过

\[
\Delta\rightarrow\Delta-\mathbf k\cdot\mathbf v
\]

进入。

接下来考察Optical molasses：从散射力得到线性摩擦。

取两束等强、相向传播的红失谐激光。对低饱和 \(s_0\ll1\)，两束光的饱和效应可以在首阶分别处理。沿 \(x\) 方向的总力为

\begin{derivation}[从双向散射力到线性摩擦]
记 Lorentz 线型
\begin{equation*}
f(\delta)=\frac{1}{1+(2\delta/\Gamma_{\rm sp})^2}.
\end{equation*}
沿 $+x$ 与 $-x$ 两束等强红失谐光，原子所受单束散射力分别为
\begin{equation*}
F_\pm(v)=\frac{\hbar k\Gamma_{\rm sp}s_0}{2}\,f(\Delta\mp kv),
\end{equation*}
其中 $\mp kv$ 来自 Doppler 失谐 $\Delta\mp kv$。总力为
\begin{equation*}
F(v)=F_+(v)-F_-(v)
=\frac{\hbar k\Gamma_{\rm sp}s_0}{2}\bigl[f(\Delta-kv)-f(\Delta+kv)\bigr].
\end{equation*}
在 $kv\ll\Gamma_{\rm sp}$ 的低速区对 $kv$ 作一阶 Taylor 展开，
\begin{equation*}
f(\Delta\mp kv)=f(\Delta)\mp kv\,f'(\Delta)+\frac{(kv)^2}{2}\,f''(\Delta)+O(v^3).
\end{equation*}
代入总力，偶次项相消而奇次项相加，
\begin{equation*}
F(v)=\frac{\hbar k\Gamma_{\rm sp}s_0}{2}\bigl[-2kv\,f'(\Delta)+O(v^3)\bigr]
=-\hbar k^2\Gamma_{\rm sp}s_0\,f'(\Delta)\,v+O(v^3).
\end{equation*}
为求 $f'(\Delta)$，把 $f$ 写成 $f(\delta)=[1+u(\delta)]^{-1}$，$u(\delta)=(2\delta/\Gamma_{\rm sp})^2$，于是
\begin{equation*}
f'(\delta)=-\frac{u'(\delta)}{[1+u(\delta)]^2},
\qquad
u'(\delta)=\frac{8\delta}{\Gamma_{\rm sp}^2},
\end{equation*}
因此
\begin{equation*}
f'(\Delta)
=-\frac{8\Delta/\Gamma_{\rm sp}^2}{\bigl[1+(2\Delta/\Gamma_{\rm sp})^2\bigr]^2}.
\end{equation*}
代回 $F(v)$ 并整理，
\begin{equation*}
F(v)=-\eta\,v+O(v^3),
\qquad
\eta=\frac{8\hbar k^2 s_0\,(\Delta/\Gamma_{\rm sp})}{\bigl[1+(2\Delta/\Gamma_{\rm sp})^2\bigr]^2}.
\end{equation*}
故阻尼系数 $\eta$ 的符号由谱线斜率 $f'(\Delta)$ 决定：红失谐 $\Delta<0$ 给出 $f'(\Delta)>0$，于是 $\eta>0$，恢复性摩擦；蓝失谐反向。即正文的 $F=-\eta v$。

\textbf{物理意义。}红失谐光频率低于原子共振，从运动原子看来迎面光子更接近共振（Doppler 频移），被更强烈吸收并自发辐射各向同性光子——净效应是减速。蓝失谐则相反，原子被加速。摩擦系数 $\eta\propto s_0\Delta/[1+(2\Delta/\Gamma)^2]^2$ 在 $\Delta=-\Gamma/2$ 处取最大值，这是激光冷却最优失谐的来源。典型铯原子 ($\lambda=852\,\si{\nm}$, $\Gamma=2\pi\times5.2\,\si{\MHz}$) 在 $s_0=0.1$ 时 $\eta\sim10^{-22}\,\si{\kg\per\s}$，对应冷却时间 $\sim10\,\si{\micro\s}$。
\end{derivation}

\[
F(v)=\frac{\hbar k\Gamma_{\rm sp}}{2}s_0
\left[
\frac{1}{1+[2(\Delta-kv)/\Gamma_{\rm sp}]^2}
-\frac{1}{1+[2(\Delta+kv)/\Gamma_{\rm sp}]^2}
\right].
\]

若

\[
|kv|\ll\Gamma_{\rm sp},|\Delta|,
\]

对速度展开到一阶，常数项因两束光对称而相消，得到

\[
\boxed{F(v)=-\eta v+O(v^3)}
\]

其中

\[
\boxed{
\eta
=-\frac{8\hbar k^2s_0(\Delta/\Gamma_{\rm sp})}
{[1+(2\Delta/\Gamma_{\rm sp})^2]^2}.
}
\]

红失谐 \(\Delta<0\) 时 \(\eta>0\)，所以速度受到线性阻尼。蓝失谐会改变符号并导致反阻尼。

接下来考察动量扩散与 Doppler 极限。

散射力虽然提供阻尼，每次自发辐射却带来随机反冲。低饱和近似下，可以把速度分布的慢演化写成 Fokker--Planck 方程

\[
\partial_t W(p,t)
=\frac{\eta}{M}\partial_p[pW]
+D_p\partial_p^2W,
\]

其中 \(D_p\) 是动量扩散系数。稳态为 Gaussian，满足 Einstein 型关系

\[
k_BT=\frac{D_p}{\eta}.
\]


本讲义采用 Fokker--Planck 形式

\[
\partial_tW=-\partial_p(AW)+D_p\partial_p^2W,
\]

所以随机过程满足

\[
\frac{\dd}{\dd t}\langle(\Delta p)^2\rangle=2D_p.
\]

在最简单的一维反冲模型中，每次散射事件包含一次定向吸收和一次随机自发辐射。两束光在 \(v=0\) 附近的总散射率为

\[
R_{\rm tot}
=\frac{\Gamma_{\rm sp}s_0}
{1+(2\Delta/\Gamma_{\rm sp})^2}
+O(s_0^2).
\]

若把自发辐射也简化成沿 \(\pm x\) 等概率反冲，则单次完整散射循环对 \(p_x\) 方差的平均贡献是 \(2(\hbar k)^2\)。Poisson 统计给出

\[
\frac{\dd}{\dd t}\langle(\Delta p_x)^2\rangle
=2(\hbar k)^2R_{\rm tot},
\]

因此

\begin{equation}
\boxed{
D_p
=(\hbar k)^2
\frac{\Gamma_{\rm sp}s_0}
{1+(2\Delta/\Gamma_{\rm sp})^2}
}
\label{eq:doppler-diffusion}
\end{equation}

在三维真实偶极辐射中，自发发射的角分布会改变数值系数，但不会改变 \(D_p\propto(\hbar k)^2R_{\rm sc}\) 的结构，也不改变 Doppler 极限来自“摩擦与反冲扩散平衡”的物理机制。

把吸收与自发辐射的反冲扩散和上一节的摩擦系数代入，在一维、低饱和二能级模型中得到

\[
\boxed{
k_BT(\Delta)
=\frac{\hbar\Gamma_{\rm sp}}{4}
\frac{1+(2\Delta/\Gamma_{\rm sp})^2}
{|2\Delta|/\Gamma_{\rm sp}}.
}
\]


确实，把 \eqnrefs{eq:doppler-diffusion} 与上一节

\[
\eta
=-\frac{8\hbar k^2s_0(\Delta/\Gamma_{\rm sp})}
{[1+(2\Delta/\Gamma_{\rm sp})^2]^2}
\]

代入 \(k_BT=D_p/\eta\)，所有 \(k^2\) 和低饱和强度因子 \(s_0\) 消去，便得到上式。温度在低饱和极限与激光强度无关，正是因为增大强度会同比增大阻尼和反冲噪声。

令 \(x=|2\Delta|/\Gamma_{\rm sp}\)，右端正比于 \(x+x^{-1}\)，所以极小值在 \(x=1\)，即

\[
\Delta=-\frac{\Gamma_{\rm sp}}{2}.
\]

由此得到 Doppler 极限

\[
\boxed{
T_D=\frac{\hbar\Gamma_{\rm sp}}{2k_B}.
}
\]

温度关于失谐的无量纲曲线以及最优红失谐位置见\figref{fig:doppler-limit}。

\begin{derivation}[从反冲 Poisson 过程到 Doppler 极限]
在时间 $\dd t$ 内发生一次散射的概率为 $R_{\rm tot}\dd t+O(\dd t^2)$，散射事件近似 Poisson 过程。记 $W(p,t)$ 为纵向动量分布，单步跳变 $p\to p+\delta p$ 的主方程为
\begin{equation*}
\partial_t W(p,t)
= R_{\rm tot}\!\int\! \dd(\delta p)\,
 \bigl[W(p-\delta p,t)-W(p,t)\bigr]\,P(\delta p),
\end{equation*}
其中 $P(\delta p)$ 为单次反冲分布。对 $W(p-\delta p,t)$ 在 $p$ 处 Taylor 展开，
\begin{equation*}
W(p-\delta p,t)
= W - \delta p\,\partial_p W
 + \frac{(\delta p)^2}{2}\,\partial_p^2 W + O\bigl((\delta p)^3\bigr).
\end{equation*}
代回主方程并按矩 $\langle\delta p\rangle=0$、$\langle(\delta p)^2\rangle=2(\hbar k)^2$ 截断到二阶（Kramers--Moyal），因此
\begin{equation*}
\partial_t W
= -\partial_p(AW)+\frac12\,\partial_p^2(BW),
\qquad
A=\langle\delta p\rangle R_{\rm tot}=0,
\quad
B=\langle(\delta p)^2\rangle R_{\rm tot}=2(\hbar k)^2 R_{\rm tot}.
\end{equation*}
扩散系数 $D_p=B/2=(\hbar k)^2 R_{\rm tot}$，即\eqnrefs{eq:doppler-diffusion}。低速摩擦 $F=-\eta v=-\eta p/M$ 给出漂移 $A(p)=-\eta p/M$，于是 Fokker--Planck 方程化为
\begin{equation*}
\partial_t W
= \frac{\eta}{M}\,\partial_p(pW)+D_p\,\partial_p^2 W.
\end{equation*}
稳态零概率流 $J=(\eta/M)pW+D_p\,\partial_p W=0$，分离变量，
\begin{equation*}
\frac{\partial_p W}{W}
= -\frac{\eta}{M D_p}\,p,
\end{equation*}
积分得
\begin{equation*}
W_{\rm ss}(p)\propto
\exp\!\left[-\frac{\eta\,p^2}{2 M D_p}\right]
=\exp\!\left[-\frac{p^2}{2 M k_B T}\right],
\qquad
k_B T=\frac{D_p}{\eta}.
\end{equation*}
代入 $D_p=(\hbar k)^2 R_{\rm tot}$ 与 $\eta=8\hbar k^2 s_0(\Delta/\Gamma_{\rm sp})/[1+(2\Delta/\Gamma_{\rm sp})^2]^2$，$k^2$ 与 $s_0$ 消去，记 $x=|2\Delta|/\Gamma_{\rm sp}$，整理得
\begin{equation*}
k_B T
= \frac{\hbar\Gamma_{\rm sp}}{4}\left(x+\frac1x\right).
\end{equation*}
由 AM--GM 不等式 $x+x^{-1}\ge2$，等号在 $x=1$ 取得，因此最优失谐 $\Delta=-\Gamma_{\rm sp}/2$（红失谐侧），于是
\begin{equation*}
k_B T_D=\frac{\hbar\Gamma_{\rm sp}}{2},
\qquad
T_D=\frac{\hbar\Gamma_{\rm sp}}{2 k_B}.
\end{equation*}

\textbf{物理意义。}Doppler 极限温度 $T_D=\hbar\Gamma/(2k_B)$ 是单频激光冷却的基本下界：冷却力来自吸收-自发辐射循环，每次自发辐射给原子反冲 $\hbar k$，这赋予加热率 $\sim\hbar^2k^2\Gamma$。冷却与加热平衡给出 $T_D$。对铯 ($\Gamma=2\pi\times5.2\,\si{\MHz}$)，$T_D\approx125\,\si{\micro\K}$；对铷 ($\Gamma=2\pi\times6.1\,\si{\MHz}$)，$T_D\approx146\,\si{\micro\K}$。突破 Doppler 极限需要 Sisyphus 冷却或蒸发冷却等额外机制，可达纳开尔文量级。
\end{derivation}

\begin{figure}[tbp]
\centering
\includegraphics[width=.83\textwidth]{doppler_limit.pdf}
\caption{低饱和二能级 Doppler 模型中的平衡温度随失谐的变化。极小值出现在\(\Delta=-\Gamma_{\rm sp}/2\)。}
\label{fig:doppler-limit}
\end{figure}


这个极限依赖``二能级、低饱和、仅 Doppler 摩擦与单光子反冲扩散''的模型。存在 Zeeman 简并和偏振梯度时，额外的亚 Doppler 机制可以把温度继续降到 \(T_D\) 以下。

\section{内部简并与亚 Doppler 冷却}
接下来考察Sisyphus 冷却的核心机制。

Doppler 极限并不是绝对下限。具有简并 Zeeman 子能级的原子在偏振梯度场中会发生位置依赖的 light shift。

例如两束正交线偏振光反向传播：

\[
\mathbf E^{(+)}(z)
=\frac{E_0}{2}
(\hat x \ee^{\ii kz}+\hat y \ee^{-\ii kz}).
\]

重新选取正交线偏振基，可看出偏振随 \(z\) 周期变化：

\[
\text{linear}\rightarrow \sigma^-\rightarrow\text{linear}\rightarrow\sigma^+\rightarrow\cdots
\]

不同 \(m_F\) 子态由于 CG 系数不同，得到不同的 AC Stark 势 \(U_{m_F}(z)\)。


以最简单的 \(J_g=1/2\to J_e=3/2\) 模型为例。对基态 \(m_g=+1/2\)，\(\sigma^+\) 跃迁到 \(m_e=+3/2\) 的 CG 强度为 1，而 \(\sigma^-\) 到 \(m_e=-1/2\) 的强度为 \(1/3\)；对 \(m_g=-1/2\) 两者交换。因此在远失谐、低饱和条件下，若局域圆偏振分量强度为 \(I_\pm(z)\)，两条基态 light shift 具有结构

\[
U_+(z)=C\left[I_+(z)+\frac13 I_-(z)\right],
\qquad
U_-(z)=C\left[I_-(z)+\frac13 I_+(z)\right],
\]

其中 \(C\propto1/\Delta\) 包含共同的偶极矩和失谐因子。lin\(\perp\)lin 场中 \(I_+\) 与 \(I_-\) 在空间上错相振荡，因此 \(U_+\) 与 \(U_-\) 也是相互错位的周期势。光抽运率又在局域圆偏振最强处最大，恰好使原子更容易在“爬到某一势能山顶”附近跳到另一条较低势曲线。Sisyphus 图像由此从“偏振变化”定量落到“CG 加权的两条位置依赖光移”。

原子沿势坡上爬，动能转化为势能；在接近势顶时被 optical pumping 到另一个更低的子态，再次开始爬坡。重复这一过程就像希腊神话中的 Sisyphus 不断推石上山，因此称 Sisyphus cooling。

接下来考察超精细结构与内部态标记。

原子电子总角动量

\[
\mathbf J=\mathbf L+\mathbf S,
\]

与核自旋 \(\mathbf I\) 耦合形成

\[
\mathbf F=\mathbf I+\mathbf J.
\]

磁偶极超精细 Hamiltonian

\[
H_{hf}=A_{\rm hf}\mathbf I\cdot\mathbf J.
\]

利用

\[
\mathbf I\cdot\mathbf J
=\frac12(\mathbf F^2-\mathbf I^2-\mathbf J^2),
\]

得到

\[
\boxed{
E_F
=\frac{A_{\rm hf}}2
[F(F+1)-I(I+1)-J(J+1)]
}.
\]

例如 \(^{133}\mathrm{Cs}\) 基态 \(6S_{1/2}\)：

\[
I=\frac72,\qquad J=\frac12,
\]

所以

\[
F=3,4.
\]

这两个超精细态之间的跃迁频率定义了 SI 秒的基础。

\section{Ramsey 干涉、频率基准与相对论修正}
接下来考察Ramsey 干涉与原子钟。

Ramsey 方法把``连续长时间驱动''改写为两个短 \(\pi/2\) 脉冲，中间留出自由相位积累时间 \(T\)。在理想瞬时脉冲近似下，取第一脉冲相位为零，\(\pi/2\) 操作可写为

\[
U_{\pi/2}=\frac1{\sqrt2}
\begin{pmatrix}
1&-\ii\\
-\ii&1
\end{pmatrix}.
\]

旋转表象中自由演化传播子为

\[
U_0(T)=
\begin{pmatrix}
 \ee^{+\ii\Delta T/2}&0\\
0&\ee^{-\ii\Delta T/2}
\end{pmatrix}.
\]

从基态 \((1,0)^T\) 出发，完整序列

\[
|\psi_f\rangle
=U_{\pi/2}U_0(T)U_{\pi/2}|g\rangle
\]

给出

\[
\boxed{
P_e(T)=\frac12[1+\cos(\Delta T)]
}
\]

（改变第二脉冲相位只会平移 fringe）。因此中央条纹的频率尺度由 \(1/T\) 决定，而不是由短脉冲宽度决定。

\begin{derivation}[Ramsey 序列作为两次 SU(2) 旋转]
记基态与激发态列向量 $|g\rangle=(1,0)^T$、$|e\rangle=(0,1)^T$。第一束同相位 $\pi/2$ 脉冲的 SU(2) 作用为
\begin{equation*}
U_{\pi/2}=\frac1{\sqrt2}
\begin{pmatrix}1&-\ii\\-\ii&1\end{pmatrix},
\end{equation*}
作用在 $|g\rangle$ 上，
\begin{equation*}
U_{\pi/2}|g\rangle
=\frac1{\sqrt2}
\begin{pmatrix}1&-\ii\\-\ii&1\end{pmatrix}
\begin{pmatrix}1\\0\end{pmatrix}
=\frac1{\sqrt2}\begin{pmatrix}1\\-\ii\end{pmatrix}.
\end{equation*}
在旋转表象中，自由演化对角传播子为
\begin{equation*}
U_0(T)=
\begin{pmatrix}\ee^{+\ii\Delta T/2}&0\\0&\ee^{-\ii\Delta T/2}\end{pmatrix},
\end{equation*}
作用后
\begin{equation*}
U_0(T)U_{\pi/2}|g\rangle
=\frac1{\sqrt2}
\begin{pmatrix}\ee^{+\ii\Delta T/2}&0\\0&\ee^{-\ii\Delta T/2}\end{pmatrix}
\begin{pmatrix}1\\-\ii\end{pmatrix}
=\frac1{\sqrt2}
\begin{pmatrix}\ee^{+\ii\Delta T/2}\\-\ii\ee^{-\ii\Delta T/2}\end{pmatrix}.
\end{equation*}
再作用第二束同相位 $\pi/2$ 脉冲，
\begin{align*}
|\psi_f\rangle
&=U_{\pi/2}\,U_0(T)\,U_{\pi/2}|g\rangle \\
&=\frac1{\sqrt2}
\begin{pmatrix}1&-\ii\\-\ii&1\end{pmatrix}
\cdot\frac1{\sqrt2}
\begin{pmatrix}\ee^{+\ii\Delta T/2}\\-\ii\ee^{-\ii\Delta T/2}\end{pmatrix} \\
&=\frac12
\begin{pmatrix}
\ee^{+\ii\Delta T/2}-\ee^{-\ii\Delta T/2}\\
-\ii\bigl(\ee^{+\ii\Delta T/2}+\ee^{-\ii\Delta T/2}\bigr)
\end{pmatrix}.
\end{align*}
由 Euler 公式 $\ee^{\ii x}+\ee^{-\ii x}=2\cos x$、$\ee^{\ii x}-\ee^{-\ii x}=2\ii\sin x$，于是
\begin{equation*}
|\psi_f\rangle
=\begin{pmatrix}
\ii\sin(\Delta T/2)\\
-\ii\cos(\Delta T/2)
\end{pmatrix},
\end{equation*}
激发态振幅
\begin{equation*}
c_e=-\ii\cos\frac{\Delta T}{2}.
\end{equation*}
取模平方，
\begin{equation*}
P_e=|c_e|^2=\cos^2\frac{\Delta T}{2}
=\frac12\bigl[1+\cos(\Delta T)\bigr].
\end{equation*}
自由演化时间 $T$ 把微小失谐放大为相位 $\Delta T$，这正是 Ramsey spectroscopy 的分辨率由 $1/T$ 而非脉冲宽度决定的来源。

\textbf{物理意义。}Ramsey 方法用两次 $\pi/2$ 脉冲中间隔自由演化时间 $T$，等效于在频域用 $\sin(\Delta T/2)/(\Delta/2)$ 窄峰探测共振。中心条纹宽度 $\sim1/T$，远窄于单脉冲的 $1/\tau_{\rm pulse}$。铯原子钟用 Ramsey 干涉在 $T\sim1\,\si{\s}$ 的喷泉时间上实现 $10^{-16}$ 频率稳定度——这是秒定义的物理基础。在量子计算中，Ramsey 序列也是刻画量子比特退相干时间 $T_2^*$ 的标准协议。
\end{derivation}

若一次测量包含 \(N\) 个互不纠缠原子，在最敏感工作点上二项分布投影噪声给出 \(\Delta P\sim1/(2\sqrt N)\)，而 \(|\partial P/\partial\phi|=1/2\)，所以

\[
\Delta\phi\sim\frac1{\sqrt N},
\qquad
\boxed{
\Delta\omega\sim\frac1{T\sqrt N}.
}
\]

这就是标准量子极限的尺度。有限脉冲时间、局域振荡器噪声、原子数涨落和 decoherence 会改变实际 Allan deviation，但不会改变这里展示的基本干涉机制。

接下来考察从频率基准到导航测距。

GPS 接收机要解四个未知量：

\[
(x,y,z,\delta t).
\]

每颗卫星给出一个伪距方程

\[
\rho_i
=c(t_{rec}-t_i)
=\sqrt{(x-x_i)^2+(y-y_i)^2+(z-z_i)^2}+c\delta t.
\]

因此至少需要四颗卫星。

时间误差 \(\delta t\) 会直接转化为空间误差

\[
\delta r\sim c\delta t.
\]

例如

\[
1\,\text{ns}\Rightarrow 0.30\,\text{m}.
\]

所以卫星导航本质上是一项高精度时间测量技术。

接下来考察重力红移。

弱场极限中，不同重力势 \(\Phi\) 下时钟频率满足

\[
\frac{\Delta\nu}{\nu}
\approx\frac{\Delta\Phi}{c^2}.
\]

在均匀重力场中，高度差 \(H\)：

\[
\Delta\Phi=gH,
\]

所以

\[
\boxed{
\frac{\Delta\nu}{\nu}\approx\frac{gH}{c^2}
}.
\]

这既可从广义相对论理解，也可借助等效原理在加速电梯中用 Doppler shift 得到。

\section{磁俘获、绝热内部态与合成规范场}
接下来考察磁阱与低场趋向态。

Zeeman 相互作用

\[
H_Z=-\boldsymbol\mu\cdot\mathbf B.
\]

绝热情况下，某个内部态的能量近似为

\[
U(\mathbf r)=\mu_{eff}|\mathbf B(\mathbf r)|.
\]

若 \(\mu_{eff}>0\)，能量在弱磁场处更低，称 low-field-seeking state，可被 \(|B|\) 的极小值束缚。

在无电流区域

\[
\nabla\cdot\mathbf B=0,
\qquad
\nabla\times\mathbf B=0,
\]

因此各分量满足 Laplace 方程。静态自由空间磁场不能产生稳定的 \(|B|\) 局域极大值，因此 high-field-seeking state 不能通过普通静态磁场极大值稳定俘获。


这个结论可以直接从 \(|\mathbf B|\) 的 Laplacian 算出来。记 \(B=|\mathbf B|\)，则

\[
\partial_iB=\frac{B_j\partial_iB_j}{B},
\]

进一步

\[
\partial_i^2B
=\frac{(\partial_iB_j)^2+B_j\partial_i^2B_j}{B}
-\frac{(B_j\partial_iB_j)^2}{B^3}.
\]

对空间指标 \(i\) 求和，并在无源区使用 \(\nabla^2B_j=0\)：

\[
\nabla^2B
=\frac{\sum_{ij}(\partial_iB_j)^2}{B}
-\frac{\sum_i(B_j\partial_iB_j)^2}{B^3}.
\]

Cauchy--Schwarz 不等式给出

\[
(B_j\partial_iB_j)^2
\le B^2\sum_j(\partial_iB_j)^2,
\]

因此

\begin{equation}
\boxed{\nabla^2|\mathbf B|\ge0}
\label{eq:magnetic-field-subharmonic}
\end{equation}

只要 \(B\neq0\)，\(|\mathbf B|\) 是 subharmonic function，因而不能在无源区有严格局域极大值。这就是静态磁阱版的 Earnshaw 限制。极小值则不被禁止，所以 low-field seeker 可以被 Ioffe--Pritchard 等非零极小值磁阱稳定束缚；保持非零偏置场还能避免零场处的 Majorana spin-flip 损失。

这不是说``超精细态不能被磁阱束缚''，相反，磁阱实验中正是选择合适的低场趋向超精细态。

接下来考察绝热基与合成规范场。

若内部 Hamiltonian 随位置变化：

\[
H_{int}(\mathbf r)|\chi_n(\mathbf r)\rangle
=\varepsilon_n(\mathbf r)|\chi_n(\mathbf r)\rangle,
\]

展开总波函数

\[
|\Psi\rangle
=\sum_n\psi_n(\mathbf r)|\chi_n(\mathbf r)\rangle.
\]

动量作用时会打到位置依赖的内部基，产生 Berry connection：

\[
\boxed{
\mathbf A_{mn}
=\ii\hbar\langle\chi_m|\nabla\chi_n\rangle
}.
\]

绝热投影到单个能带后

\[
H_{eff}
=\frac1{2M}[\mathbf p-\mathbf A_n(\mathbf r)]^2
+\varepsilon_n(\mathbf r)+\Phi_n(\mathbf r).
\]


其中单带几何标量势来自被绝热消去的其他内部态：

\[
\boxed{
\Phi_n(\mathbf r)
=\frac{\hbar^2}{2M}
\sum_{m\neq n}
|\langle\chi_m|\nabla\chi_n\rangle|^2.
}
\]

Berry curvature 则为

\[
\mathbf B_{\rm syn}=\nabla\times\mathbf A_n.
\]

因此“合成磁场”并不是把真实磁场直接替换成另一个实验旋钮，而是原子质心在位置依赖内部本征基中绝热运动时积累的几何相位所产生的有效 Lorentz 结构。这与\chapref{ch:born-huang}中的分子核几何规范势是同一个数学机制。

这就是冷原子中 synthetic gauge field 的基本来源。

\begin{derivation}[lin\(\perp\)lin 构型中 Sisyphus 冷却的摩擦系数与极限温度，即\eqnrefs{eq:sisyphus-temperature}]
\emph{策略路线图。}首先写出 lin\(\perp\)lin 光场中两个基态子能级的光位移势与局域光抽运率；其次在准经典轨道近似下计算原子穿越一个偏振周期时的平均能量耗散；再次由能量耗散率提取速度线性摩擦系数 \(\alpha_{\rm Sis}\)；然后计算自发辐射反冲导致的动量扩散系数 \(D_p\)；最后由 Einstein 关系 \(k_BT=D_p/\alpha_{\rm Sis}\) 得到 Sisyphus 极限温度，并证明它远低于 Doppler 极限。

\emph{第一步：光位移势与抽运率。}在 \(J_g=1/2\to J_e=3/2\) 跃迁的 lin\(\perp\)lin 构型中，两束等强反向传播的线偏振光产生空间周期偏振梯度。记驻波强度为 \(I_0\)，饱和参数 \(s_0=I_0/I_{\rm sat}\)，远失谐 \(|\Delta|\gg\Gamma\)。两个基态 \(m_g=\pm1/2\) 的 AC Stark 光位移为
\begin{equation}
U_\pm(z)=\frac{U_0}{2}\bigl[1\pm\cos(2kz)\bigr],
\label{eq:sisyphus-light-shift}
\end{equation}
其中 \(U_0=\hbar\Omega^2/(2\Delta)\) 是光位移幅度，\(\Omega\) 为 Rabi 频率。局域光抽运率（从高势态到低势态的自发跃迁率）为
\begin{equation}
\Gamma_{\rm pump}(z)=\Gamma'\sin^2(2kz),
\qquad
\Gamma'=\frac{\Gamma s_0}{2}\frac{1}{1+(2\Delta/\Gamma)^2}\approx\frac{\Gamma\Omega^2}{4\Delta^2},
\label{eq:sisyphus-pump-rate}
\end{equation}
远失谐近似下最后一步用了 \(s_0=2\Omega^2/\Gamma^2\)。关键特征：\(U_+\) 的势顶（\(kz=\pi/2\)）恰是 \(\Gamma_{\rm pump}\) 的极大处——原子在势顶最容易被抽运到另一条势谷。

\emph{第二步：准经典轨道与单周期能量耗散。}考虑原子以速度 \(v\) 沿 \(z\) 方向运动，\(z=vt\)。在一个偏振半周期 \(\lambda/2=\pi/k\) 内，原子从 \(U_+\) 势谷出发爬到势顶。在准经典近似 \(E_{\rm kin}\gg|U_0|\) 下，原子位置随时间近似线性。势顶附近 \(\cos(2kz)\approx-1+2[kz-\pi/2]^2\)，抽运率展开为
\begin{equation*}
\Gamma_{\rm pump}(z)\approx\Gamma'\bigl[1-4(kz-\pi/2)^2+\cdots\bigr].
\end{equation*}
原子穿越势顶邻域 \(|kz-\pi/2|<\delta\) 的时间为 \(\Delta t=2\delta/(kv)\)。在此区间内未被抽运的概率为
\begin{equation*}
P_{\rm survive}=\exp\!\Bigl[-\int\Gamma_{\rm pump}\,\dd t\Bigr]
\approx\exp\!\Bigl[-\frac{2\Gamma'}{kv}\int_0^\delta(1-4u^2)\,\dd u\Bigr]
=\exp\!\Bigl[-\frac{2\Gamma'}{kv}\bigl(\delta-\tfrac{4}{3}\delta^3\bigr)\Bigr].
\end{equation*}
取 \(\delta\to\infty\)（抽运主要发生在势顶附近，远端贡献指数小），单周期内被抽运的概率
\begin{equation*}
P_{\rm pump}^{(\rm per)}\approx 1-\exp\!\bigl[-\pi\Gamma'/(kv)\bigr].
\end{equation*}
每次抽运从 \(U_+\) 势顶（\(U_+\approx U_0\)）跳到 \(U_-\) 势谷（\(U_-\approx 0\)），平均能量损失
\begin{equation*}
\langle\Delta E\rangle_{\rm per}\approx -U_0\,P_{\rm pump}^{(\rm per)}.
\end{equation*}

\emph{第三步：摩擦系数。}在低速区 \(kv\ll\Gamma'\)（即原子在穿越一个周期内有大概率被抽运），展开 \(P_{\rm pump}^{(\rm per)}\approx\pi\Gamma'/(kv)\)。单位时间穿越周期数为 \(2v/\lambda=kv/\pi\)，故平均能量耗散率
\begin{equation*}
\frac{\dd E}{\dd t}
=\frac{kv}{\pi}\langle\Delta E\rangle_{\rm per}
\approx-\frac{kv}{\pi}\,U_0\,\frac{\pi\Gamma'}{kv}
=-U_0\Gamma'.
\end{equation*}
此结果与速度无关——在低速区 Sisyphus 冷却给出常功率阻尼而非线性摩擦。为得到线性摩擦区，需考虑 \(E_{\rm kin}\sim|U_0|\) 的量子区：原子被囚禁于单个势阱内，能级间距 \(\hbar\omega_{\rm opt}\sim\sqrt{8U_0E_r}\)（\(E_r=\hbar^2k^2/2M\) 为反冲能）。振动频率 \(\omega_{\rm opt}=2k\sqrt{|U_0|/2M}\)。每次振动周期损失能量 \(\sim|U_0|\)，阻尼率 \(\gamma_E\sim\Gamma'U_0/(\hbar\omega_{\rm opt})\)。对应速度阻尼系数
\begin{equation}
\alpha_{\rm Sis}=\frac{M\gamma_E}{2}
\sim\frac{\Gamma'|U_0|M}{2\hbar\omega_{\rm opt}}
=\frac{\Gamma'M\sqrt{|U_0|}}{2\hbar k\sqrt{8M}}
\propto\frac{\hbar k^2\Gamma'\sqrt{|U_0|}}{\sqrt{E_r}}.
\label{eq:sisyphus-friction}
\end{equation}

\emph{第四步：动量扩散与极限温度。}动量扩散来自自发辐射随机反冲。抽运率 \(\Gamma'\) 每次发射反冲 \(\hbar k\)，扩散系数
\begin{equation*}
D_p\sim(\hbar k)^2\Gamma'\times\frac{1}{2}
=\frac{\hbar^2k^2\Gamma'}{2}.
\end{equation*}
由 Einstein 关系 \(k_BT=D_p/\alpha_{\rm Sis}\)，取 \(\alpha_{\rm Sis}\) 在量子区的表达式并注意 \(\omega_{\rm opt}\) 与 \(U_0\) 的关系，
\begin{equation}
k_BT_{\rm Sis}\sim\frac{D_p}{\alpha_{\rm Sis}}
\sim\frac{\hbar^2k^2\Gamma'/2}{\Gamma'M\sqrt{|U_0|}/(2\hbar k\sqrt{8M})}
=\frac{\hbar\omega_{\rm opt}}{4}\sim\frac{|U_0|}{4},
\label{eq:sisyphus-temperature}
\end{equation}
即 Sisyphus 极限温度正比于光位移幅度 \(|U_0|\)。由于远失谐下 \(|U_0|=\hbar\Omega^2/(2|\Delta|)\) 可通过减小光强或增大失谐任意降低，而 Doppler 极限 \(T_D=\hbar\Gamma/(2k_B)\) 固定，故
\begin{equation*}
\frac{T_{\rm Sis}}{T_D}\sim\frac{|U_0|/4}{\hbar\Gamma/2}
=\frac{|U_0|}{2\hbar\Gamma}
=\frac{\Omega^2}{4|\Delta|\Gamma}
=\frac{s_0\Gamma}{8|\Delta|}\ll1
\end{equation*}
在远失谐 \(|\Delta|\gg\Gamma\) 且低饱和 \(s_0\ll1\) 时严格成立。即\eqnrefs{eq:sisyphus-temperature}。

\emph{物理意义。}Doppler 极限来自散射力冷却与自发辐射反冲加热的平衡，两者都正比于 \(\Gamma\)，故 \(T_D\propto\hbar\Gamma\)。Sisyphus 冷却的冷却力来自光位移势的有序能量转化（\(\propto|U_0|\)），而加热仍来自自发辐射反冲（\(\propto\hbar^2k^2\Gamma'\)），两者参数独立可调，故极限温度可远低于 Doppler 极限。实验上铯原子 Sisyphus 冷却可达 \(T\sim2\,\si{\micro\K}\)，比 \(T_D\approx125\,\si{\micro\K}\) 低约 60 倍。


\emph{[Doppler 冷却极限与反冲极限的精确推导]} 激光冷却的两个基本极限——Doppler 极限和反冲极限——由不同的物理机制决定，其精确推导揭示冷却的量子极限。分四步建立。

\emph{第一步：Doppler 冷却的辐射压力。}二能级原子在红失谐激光（\(\delta=\omega_L-\omega_0<0\)）中，辐射压力 \(F=\hbar k\cdot\Gamma_{\rm sc}\)（\(\Gamma_{\rm sc}\) 散射率）依赖原子速度 \(v\) 通过 Doppler 频移 \(\delta_{\rm eff}=\delta-kv\)。低速展开 \(F\approx-\alpha v\)（\(\alpha=-\hbar k^2\frac{\partial\Gamma_{\rm sc}}{\partial\delta}\big|_{\delta}\)），\(\alpha>0\) 给出粘滞阻尼。同时散射引入动量扩散 \(D_p=\hbar^2k^2\Gamma_{\rm sc}/2\)（每散射事件注入 \(\hbar k\) 随机动量）。

\emph{第二步：Doppler 极限。}平衡温度由粘滞-扩散平衡 \(k_BT=D_p/\alpha\) 给出。优化失谐 \(\delta=-\Gamma/2\)（\(\Gamma\) 自然线宽）给出 Doppler 极限
\begin{equation}
  k_BT_D=\frac{\hbar\Gamma}{2},
  \label{eq:phys-lc-doppler-limit}
\end{equation}
即 Doppler 极限温度 \(T_D=\hbar\Gamma/(2k_B)\)。对铯 \(T_D\approx125\,\mu\)K，对铷 \(T_D\approx146\,\mu\)K。Doppler 极限来自辐射压力的涨落（自发辐射的随机性），是二能级近似的根本限制。

\emph{第三步：Sisyphus 冷却与反冲极限。}多能级原子（含 Zeeman 子能级）在偏振梯度光中经历 Sisyphus 冷却：光位移 \(U(\boldsymbol r)\) 依赖偏振（\(\sigma^+\)/\(\sigma^-\) 分量空间变化），原子在势丘上爬时通过光学泵浦"翻转"到谷底——反复爬丘但总在谷底，如 Sisyphus 推石。Sisyphus 冷却突破 Doppler 极限，达到反冲极限
\begin{equation*}
  k_BT_R=\frac{\hbar^2k^2}{2m},
\end{equation*}
即单光子反冲能。对铯 \(T_R\approx2\,\mu\)K，远低于 \(T_D\)。Sisyphus 冷却的关键是光位移的内部态依赖（多能级效应），二能级近似无法实现。

\emph{第四步：亚反冲冷却与 VSCPT。}速度选择相干布居阱（VSCPT）冷却到亚反冲温度 \(T<T_R\)：利用暗态（速度 \(v=0\) 的相干叠加态不散射光子）将原子聚集到 \(v=0\)。VSCPT 的平衡温度由冷却时间 \(t\) 决定 \(T\sim\hbar/(kt)\)（随时间无下限降温，实际由加热机制限制）。拉曼冷却和蒸发冷却可进一步降至纳开尔文——这是原子钟和 BEC 实验的温度范围。冷却极限的层级 \(T_D\to T_R\to T_{\rm sub-recoil}\) 反映了量子极限的逐步突破：Doppler（涨落）→ 反冲（量子离散）→ 亚反冲（相干暗态），每步利用更多量子相干性。
\end{derivation}
